% =====================================================================
%  MINI THESIS PREVIEW  —  for review BEFORE editing the real Overleaf thesis
%  Written 2026-06-10. Compiles standalone (pdflatex, no bib needed).
%
%  PURPOSE: let the author read the RESTRUCTURED story + structure (the
%  check.txt spine, A+D narrative, IMRAD) with the CURRENT best numbers,
%  approve it, THEN expand into docs/thesis_working_distilling_overleaf/
%  and swap PROVISIONAL metrics for the Tier-2 full-frame run.
%
%  METRIC STATUS LEGEND (every result is tagged):
%    [FINAL]       at-standard, recorded in knowledge/evals.csv -> will NOT
%                  change with the full run.
%    [PROVISIONAL] from the n=1000 mini-run or pending the 4k strided cache
%                  -> WILL be replaced by Tier-1 (4k) then Tier-2 (full).
%  Each result sentence carries a % [source: eval=<id>] comment for tracing.
% =====================================================================
\documentclass[11pt]{article}
\usepackage[a4paper,margin=1in]{geometry}
\usepackage{booktabs}
\usepackage{amsmath}
\usepackage{xcolor}
\usepackage{enumitem}
\usepackage{hyperref}
\hypersetup{colorlinks=true,linkcolor=blue!50!black,urlcolor=blue!50!black}
\newcommand{\final}{{\small\textcolor{green!45!black}{[FINAL]}}}
\newcommand{\prov}{{\small\textcolor{orange!80!black}{[PROVISIONAL]}}}
\newcommand{\todo}[1]{{\small\textcolor{red!70!black}{[\,#1\,]}}}
\title{\textbf{Cross-Modal Drone Detection:\\
A Self-Improving, Deployable Detection Pipeline}\\[4pt]
{\large\itshape Mini preview --- structure \& story for review (not the final draft)}}
\author{Ahmed}
\date{2026-06-10}

\begin{document}
\maketitle

\begin{abstract}
\noindent This document is a \emph{compressed preview} of the restructured thesis:
the narrative spine, the IMRAD skeleton, and the empirical tables, with current-best
numbers tagged \final{} (locked) or \prov{} (to be replaced by the full-frame run).
Read it to approve the \emph{shape of the argument}; the full chapters are expanded
only after sign-off.

\medskip
\noindent\textbf{Thesis in one paragraph.} We build and validate a \emph{deployable}
dual-modality (RGB + thermal-IR) drone-detection \emph{system} --- a fusion pipeline with
an operator GUI --- and we show it is produced and \emph{maintained} by a self-improving
loop: a human-in-the-loop (HITL) data engine that co-develops the detectors, and a
statistics-first feature-space diagnostic (``Model MRI'') that derives cheap, robust
components and catches regressions before they ship. The headline is the \emph{system and
the method that sustains it} (contributions C1--C3); cross-modal grayscale transfer, a
28-point scoring-rule audit, and large speed wins are supporting \emph{findings}.
\end{abstract}

% ---------------------------------------------------------------
\section{Introduction}
\label{sec:intro}

\paragraph{Problem.} Drone detection is reported on saturated academic benchmarks, but a
detector that wins a benchmark is not a system that survives deployment. Real operation adds
(i) confusers --- birds, aircraft, helicopters --- that a recall-tuned detector hallucinates as
drones; (ii) hard surfaces (low-resolution CCTV, distant/small targets) where benchmark recall
collapses; and (iii) degraded sensing (a thermal channel that is really a desaturated RGB feed).
The gap between \emph{benchmark score} and \emph{deployable system} is the problem this thesis
addresses.

\paragraph{Research questions.}
\begin{description}[leftmargin=2.2em,itemsep=2pt]
  \item[RQ1] Can a learned, leakage-aware pipeline suppress out-of-distribution confuser
        false positives without sacrificing drone recall?
  \item[RQ2] What is the in-distribution cost of that pipeline, and how does it depend on the
        surface (clean benchmark vs.\ hard CCTV)?
  \item[RQ3] Does a thermal-trained detector transfer to a degraded \emph{grayscale} channel,
        and can the same diagnostic make that transfer reliable?
\end{description}

\paragraph{Contributions.}
\begin{description}[leftmargin=2.2em,itemsep=2pt]
  \item[C1 --- The pipeline + GUI.] A deployable dual-modality detection pipeline (per-modality
        detect $\rightarrow$ trust-routing classifier $\rightarrow$ feature-space verifier
        $\rightarrow$ temporal alert-gate) with an operator GUI. The \emph{artifact}.
  \item[C2 --- The Model MRI.] A statistics-before-training, feature-space diagnostic
        (LDA/PCA/ANOVA/AUROC/leakage/calibration) that tells us \emph{what will work before we
        train}. It auto-derives the cheap MLP verifier and the leakage-aware trust router, and
        it catches regressions.
  \item[C3 --- The HITL data engine.] A human-in-the-loop label-reviewer that co-develops the
        detectors and provides the discipline that catches a real regression (the IR V5 event).
\end{description}
Supporting \emph{findings} (not the spine): cross-modal grayscale transfer (RQ3); a scoring-rule
audit; and 37--404$\times$ component speed-ups.

% ---------------------------------------------------------------
\section{Background \& Related Work}
\label{sec:background}
\todo{expand later}: benchmark saturation (Anti-UAV); confuser robustness; cross-modal transfer;
HITL annotation; feature-space diagnostics. One paragraph each, positioning C1--C3 against prior work.

% ---------------------------------------------------------------
\section{Methodology}
\label{sec:method}

\subsection{The system (C1)}
The pipeline scores each modality independently and never unions their evidence. Stages:
(1) \textbf{detectors} --- RGB \texttt{ft4} and thermal-IR \texttt{v3b}; (2) a \textbf{trust-routing
classifier} (\texttt{robust8}) that decides, per frame, which modality to believe
(trust-RGB / trust-IR / trust-both / reject); (3) a per-detection \textbf{verifier}
(\texttt{mlp\_v5} for RGB, \texttt{mlp\_v5\_ir\_aligned} for thermal \emph{and} grayscale) that
drops confuser detections using the detector's own ROI features; (4) a \textbf{temporal alert-gate}
(segment-level 2-of-3 voting, mirrored from the GUI). \textbf{GUI}: \todo{real screenshot needed} ---
operator view, per-modality tracks, alert gate, grayscale fallback path.

\subsection{The HITL data engine (C3)}
Detectors are co-developed with a label-reviewer in the loop: every mined batch is human-vetted
before it enters training. Section~\ref{sec:partA} shows what happens when that discipline is
\emph{bypassed} (the IR V5 regression).

\subsection{The Model MRI (C2)}
Before any training run we diagnose separability in feature space. The detector's own ROI features
\emph{already} separate drones from confusers --- LDA accuracy
$0.952$ (RGB) / $0.981$ (IR)~\final{} % [source: eval=mri_lda]
and a dominant ANOVA effect ($F=42{,}346$)~\final{} % [source: eval=mri_anova]
--- which is the hypothesis that justifies a \emph{cheap MLP} verifier instead of re-processing
every crop with a CNN. The same diagnostic shows the grayscale channel is \emph{alignable} to the
thermal feature space: per-modality $z$-alignment lifts held-out gray$\rightarrow$thermal
discrimination from AUROC $0.500$ to $0.919$~\final{}. % [source: eval=mri_gray_align]
\textbf{This is the through-line:} a hypothesis in feature space $\rightarrow$ figures $\rightarrow$
a derived component $\rightarrow$ a pipeline win.

\subsection{Component design as a three-version argument}
Each component is presented as three versions --- a \emph{sad} path, a \emph{happy} path, and the
step that \emph{justifies the next thing}:
\begin{itemize}[leftmargin=1.4em,itemsep=2pt]
  \item \textbf{RGB detector}: baseline / retrained\_v2 (OOD recall collapse) / \texttt{ft4}.
        The collapse cannot be fixed at the detector $\rightarrow$ motivates the verifier.
  \item \textbf{IR detector}: V2 / V5 (HITL regression) / \texttt{v3b}. Motivates HITL discipline (C3).
  \item \textbf{Trust router}: \texttt{sa32} + \texttt{fusion\_no\_fn} (32--40 feat, expensive, over-rejects)
        / \texttt{robust8} (8 feat). Motivates leakage-aware feature selection (C2).
  \item \textbf{Verifier}: MobileNetV3 patch CNN (slow, OOD-weak) / \texttt{mlp\_v5} (distilled).
        Motivates the MRI: are the features separable? Yes (LDA $\approx0.95$) $\rightarrow$ distilled MLP wins.
\end{itemize}

\subsection{Evaluation design}
\textbf{Surfaces (7 + grayscale):} RGB test, IR test, Anti-UAV (paired), Svanstr\"om (paired),
selcom CCTV, RGB confusers, IR confusers, and a grayscale path. Birds are a \emph{category}
inside the confuser sets, not a dataset.
\textbf{Image size:} $640$ default; $1280$ only for low-native-resolution sources (Svanstr\"om
$640{\times}512$, selcom CCTV).
\textbf{Scoring is per-modality, never unioned:} trust-RGB scores RGB detections vs.\ RGB GT only
(IR ground truth is \emph{excluded}, not counted as a miss); only reject-both counts both as misses.
Svanstr\"om uses IoP@0.5, Anti-UAV IoU@0.5; confusers have no GT (false positives / fire-rate only).
\textbf{Two tiers:} Tier-1 is a $\sim$4k \emph{strided} screen (per-frame numbers, this preview);
Tier-2 is the full consecutive-frame run (final numbers + temporal). Temporal voting needs
consecutive frames, so it is a placeholder until Tier-2.

% ---------------------------------------------------------------
\section{Empirical Evaluation}
\label{sec:empirical}

\paragraph{Metric status.} Part~A (detectors), the IR arc, the MRI numbers, and speed are
\final{} (at-standard, recorded). Part~B/C/D pipeline cells are \prov{} and will be replaced by
the 4k strided screen and then the full-frame run; the structure does not change, only the digits.

% ---------------- PART A -----------------
\subsection{Part A --- Per-model, in-domain (bare detectors)}
\label{sec:partA}

\begin{table}[h]\centering
\caption{RGB detector, three versions. In-domain all three are competitive; the differences are
\emph{out-of-distribution}. \texttt{ft4}'s low \emph{bare} Svanstr\"om F1 is precision traded for
recall (R$=0.914$) --- the pipeline recovers it (Part~B). \final{} except where noted.}
\begin{tabular}{lccc}
\toprule
surface (imgsz, rule) & baseline & retrained\_v2 & \texttt{ft4} (prod) \\
\midrule
rgb\_dataset\_test (640, IoU) & 0.942 & \textbf{0.949} & 0.929 \\ % [source: eval=rgb_rgbds_baseline; rgb_rgbds_retrainedv2; v5_rgbds_bare]
Svanstr\"om (1280, IoP) & \textbf{0.950} & 0.462 & 0.596\,(bare) \\ % [source: eval=rgb_svan_baseline; rgb_svan_retrainedv2; v5_svan_bare]
selcom CCTV (1280, IoP) & 0.145 & 0.007 & \textbf{0.591} \\ % [source: eval=selcom_baseline_1280; selcom_retrainedv2_1280; v5_selcom_bare]
Anti-UAV (saturated) & 0.994 & 0.994 & 0.986 \\ % [source: eval=rgb_antiuav_baseline; rgb_antiuav_retrainedv2; v5_antiuav_bare]
\bottomrule
\end{tabular}
\end{table}

\noindent\textbf{Reading.} There is \emph{no single best bare detector}: baseline wins Svanstr\"om,
\texttt{ft4} is the only one that holds CCTV ($0.591$ vs.\ $0.145$), and retrained\_v2 wins nothing
out-of-domain (Svanstr\"om $0.462$, selcom $0.007$ --- a single true positive in $295$). Crucially,
retrained\_v2 is \emph{fine in-domain} ($0.949$): a model chosen on the own-split benchmark is the
one that fails in deployment. \final{} % [source: ledger=rgb-collapse-ood-specific]
This is the entire motivation for the pipeline (C1): the best \emph{deployable} configuration is
not the best bare detector, it is \texttt{ft4}\,+\,pipeline.

\begin{table}[h]\centering
\caption{IR detector evolution (own test split, 640, IoU). V5 is a HITL regression: a large
positive batch bypassed per-batch review, so unlabelled drones scored as false positives. The
HITL discipline (C3) is what recovered it. \final{}}
\begin{tabular}{lcccl}
\toprule
version & P & R & F1 & \\
\midrule
V2  & 0.661 & 0.406 & 0.503 & early, weak recall \\ % [source: eval=ir_final_v2]
V5  & --    & --    & 0.737 & HITL regression \\       % [source: eval=ir_final_v5]
\texttt{v3b} (prod) & 0.957 & 0.977 & \textbf{0.967} & corrective FT \\ % [source: eval=ir_final_v3b]
\bottomrule
\end{tabular}
\end{table}

% ---------------- PART B -----------------
\subsection{Part B --- Full-pipeline ablation (paired surfaces)}
\label{sec:partB}
On the paired drone-positive surfaces (Anti-UAV, Svanstr\"om) we ablate
bare $\rightarrow$ +router $\rightarrow$ +verifier $\rightarrow$ full cascade, scored per-modality.
The question is whether precision rises \emph{without} losing recall.

\begin{table}[h]\centering
\caption{Svanstr\"om pipeline (IoP@0.5). The verifier recovers \texttt{ft4}'s bare precision deficit.
\prov{} (mini-run / at-standard head-to-head; refresh from 4k then full).}
\begin{tabular}{lccc}
\toprule
stage & P & R & F1 \\
\midrule
\texttt{ft4} bare        & 0.443 & 0.914 & 0.596 \\ % [source: eval=v5_svan_bare]
+ patch verifier (MobileNet) & 0.686 & 0.872 & 0.768 \\ % [source: eval=v5_svan_patch]
+ \texttt{mlp\_v5} verifier   & \textbf{0.903} & 0.838 & \textbf{0.869} \\ % [source: eval=v5_svan_mlp]
\bottomrule
\end{tabular}
\end{table}

\noindent Recall is preserved (R $0.914\!\rightarrow\!0.838$) while precision roughly doubles
($0.443\!\rightarrow\!0.903$). The trust router (\texttt{robust8}) adds modality selection on top;
mini-run thermal F1 $=0.979$~\prov{}. % [source: eval=robust8-full-pipeline-cmp]
\todo{4k}: fill the full bare/+C/+F/full grid for Anti-UAV and Svanstr\"om from the unified cache.

% ---------------- PART C -----------------
\subsection{Part C --- Confuser false-positive reduction (OOD benefit)}
\label{sec:partC}
On confuser surfaces (no drones present, so every surviving detection is a false positive) the
verifier is the benefit. \texttt{mlp\_v5} cuts RGB-confuser false positives $216\!\rightarrow\!16$
($\sim$13$\times$)~\prov{} % [source: eval=robust8-full-pipeline-cmp]
and, in the thermal/grayscale-aligned feature space, held-out CBAM F1 rises
$0.699\!\rightarrow\!0.841$ with false positives $48\!\rightarrow\!13$~\final{}. % [source: eval=ir-grayscale-harvest-solves-thermal-verifier]
The MobileNet patch CNN achieves comparable suppression but $37$--$72\times$ slower (Section~\ref{sec:speed}).
\todo{4k}: rgb / ir / gray confuser FP grid (bare $\rightarrow$ +mlp $\rightarrow$ +patch) from the cache.

% ---------------- PART D -----------------
\subsection{Part D --- Grayscale cross-modal transfer (finding, RQ3)}
\label{sec:partD}
We report \emph{only} the working configuration: the thermal detector \texttt{v3b} run on a
\emph{grayscale} RGB feed, plus the aligned verifier, with the (thermal-trained) router bypassed.
On held-out grayscale confusers the aligned verifier cuts false positives
$325\!\rightarrow\!12$ ($96\%$) while keeping recall within $\sim$5 points of bare~\final{}, % [source: eval=ir_aligned_gray_heldout]
and the MRI alignment that makes this possible lifts gray$\rightarrow$thermal AUROC
$0.500\!\rightarrow\!0.919$ (Section~\ref{sec:method}). The degenerate configurations (router on
grayscale; a grayscale-specific verifier that over-vetoes) are mentioned in one honest sentence
and not tabulated.

% ---------------- SPEED -----------------
\subsection{Speed}
\label{sec:speed}
\begin{table}[h]\centering
\caption{Component cost (the cheap-by-design payoff of C2). \final{} (bench).}
\begin{tabular}{lccc}
\toprule
component & sad & happy & speed-up \\
\midrule
trust router  & \texttt{fusion\_no\_fn} 38.3 ms/frame & \texttt{robust8} 0.095 ms/frame & $\sim$404$\times$ \\
verifier      & patch 59--112 ms/det & \texttt{mlp\_v5} 1.3--2.1 ms/det & $\sim$37--72$\times$ \\
\bottomrule
\end{tabular}
\end{table}
\noindent Total pipeline overhead is $\sim$1--4\%. \todo{wire bench numbers to knowledge/ledger before final}.

% ---------------- THREATS -----------------
\subsection{Threats to validity}
\label{sec:threats}
\begin{itemize}[leftmargin=1.4em,itemsep=2pt]
  \item \textbf{Strided screen.} Tier-1 numbers are a strided sample; temporal continuity is not
        yet measured. Tier-2 (full consecutive frames) resolves both.
  \item \textbf{Saturated benchmark.} Anti-UAV is saturated; it is a sanity floor, not a discriminator.
  \item \textbf{Per-modality scoring choice.} Excluding the untrusted modality's GT is a deliberate
        rule; we report and justify it rather than silently unioning (the 28-point audit).
  \item \textbf{Train/eval provenance.} Components are scored by pipeline \emph{effect}, not their own
        held-out accuracy, because those splits are not cross-comparable.
\end{itemize}

% ---------------------------------------------------------------
\section{Conclusion}
A benchmark detector is not a deployable system. We close that gap with a per-modality pipeline and
operator GUI (C1), kept honest and cheap by a feature-space diagnostic (C2) and a human-in-the-loop
data engine (C3). The components are not individually dominant --- no single bare detector wins every
surface --- but their \emph{composition} is the only configuration deployable across clean benchmarks,
hard CCTV, confuser-rich scenes, and a degraded grayscale channel.

\vspace{1em}
\hrule
\vspace{4pt}
\noindent{\small\textbf{Review checklist for the author.}
(1) Is the A+D spine (system as headline, grayscale as finding) what you want?\;
(2) Is the Part~A$\rightarrow$B framing of the ``no single best detector'' tension acceptable?\;
(3) Sign-off gates the full expansion into the Overleaf thesis; Tier-2 then swaps every \prov{} digit.}

\end{document}
